\documentclass[11pt,a4paper]{article}

\usepackage[utf8]{inputenc}
\usepackage[T1]{fontenc}
\usepackage{amsmath,amssymb,amsfonts,amsthm}
\usepackage{geometry}
\usepackage{booktabs}
\usepackage{hyperref}

\geometry{margin=2.35cm}

\newtheorem{assumption}{Assumption}
\newtheorem{theorem}{Theorem}
\newtheorem{lemma}{Lemma}
\newtheorem{proposition}{Proposition}
\newtheorem{corollary}{Corollary}
\newtheorem{remark}{Remark}

\newcommand{\R}{\mathbb R}
\newcommand{\E}{\mathbb E}
\newcommand{\op}{\mathrm{op}}
\newcommand{\fro}{\mathrm F}
\newcommand{\Id}{I}
\newcommand{\wh}{\widehat}
\newcommand{\calC}{\mathcal C}
\newcommand{\calE}{\mathcal E}
\newcommand{\tr}{\operatorname{tr}}
\newcommand{\diag}{\operatorname{diag}}
\newcommand{\Span}{\operatorname{span}}

\title{Replica Bounds for the PDE Encoder and Prompt-Conditioned Polynomial Decoder}
\author{}
\date{}

\begin{document}
\maketitle

\begin{remark}[Scope of this manuscript]
The first decoder calculation below concerns the historical collapsed
one-shot map \(W_{\rm dec}(K/M)\widehat G^\top\widehat b\).  The later
sections extend its spectral order parameters to the prompt-conditioned
matrix-free preconditioner and give the exact fixed-polynomial substitution
for recurrent Heavy--Ball and Chebyshev.  The rigorous controller-dependent
transfer audit is given in
\texttt{experiments/transformers/three\_controller\_encoder\_decoder\_generalization.tex}.
Closing the full PDE asymptotic still requires the matrix-valued Ritz moments
and teacher/physical overlaps stated below.  PCG additionally requires a
layer-indexed state evolution because its Krylov coefficients depend on the
right-hand side.  The gradient-flow ODEs in this manuscript do not describe
Muon updates.
\end{remark}

\section{Claim}

The exact task-level correspondence with the low-rank ICL regression framework
of \cite{takanami2026lowrankicl} is:
\[
    \text{one task} \quad\Longleftrightarrow\quad
    \text{one coefficient vector } z,
\]
and one prompt for that task is the whole collection
\[
    \calC_z=\{(f_i,u_i)\}_{i=1}^{m}.
\]
The weak rows inside the prompt are examples for the same task; they are not
independent tasks.  There are two uses of the same replica framework.

\begin{enumerate}
    \item \textbf{Encoder ICL.}  The weak-form encoder is trained so that the
    produced system is consistent:
    \[
        \wh Gz-\wh b\approx 0 .
    \]
    This is a low-rank regression generalization error at the level of a full
    prompt for a fixed $z$.

    \item \textbf{Decoder ICL.}  The decoder is taken to be \emph{linear}
    attention, not softmax attention.  Linear attention implements
    \[
        \wh z
        =
        W_{\rm dec}
        \left(\frac{K}{M}\wh G^\top\wh b\right),
        \qquad M=mR,
    \]
    so $W_{\rm dec}$ is a learned spectral preconditioner for the empirical
    Gram matrix.  The paper's replica solution is precisely a distributional
    formula for this learned preconditioner.
\end{enumerate}

The resulting end-to-end bound has the exact form
\[
    \boxed{
    \frac1K\E\|\wh z-z\|_2^2
    \;\lesssim\;
    E_{\rm dec}^{\rm rep,GP}
    +
    B_{\rm dec}^2 C_{\rm GP}
    E_{\rm enc}^{\rm rep,GP}
    }
\]
and
\[
    \boxed{
    \E\frac{\|\wh u_\star-u_\star\|_2^2}{\|u_\star\|_2^2}
    \;\lesssim\;
    C_{\rm elliptic}^2 B_\Psi^2K
    \left(
    E_{\rm dec}^{\rm rep,GP}
    +
    B_{\rm dec}^2 C_{\rm GP}
    E_{\rm enc}^{\rm rep,GP}
    \right)
    +\text{dictionary/gauge bias}.
    }
\]
Here $E_{\rm enc}^{\rm rep,GP}$ and $E_{\rm dec}^{\rm rep,GP}$ denote the
replica errors for the exact GP-induced prompt law.  They reduce to the closed
forms in the paper, with $D$ replaced by $K$, only in the special case where
the weak-row design is task-independent and isotropic.

\section{Exact Gaussian PDE Prompt Model}

We work with the discretized elliptic family
\[
    A(z)=A_0+\sum_{k=1}^{K} z_kA_k,
    \qquad
    A(z)u=f.
\]
Assume a Gaussian low-rank latent law
\[
    z=\sqrt{\frac{K}{r_z}}\,B\eta,
    \qquad
    \eta\sim N(0,I_{r_z}),
    \qquad
    B^\top B=I_{r_z},
\]
so the task rank is $r_z$ and $\rho_z=r_z/K$.  The full-rank Gaussian case is
included by taking $r_z=K$.

Assume the forcing is a Gaussian process.  After discretization,
\[
    f_i\sim N(0,\Sigma_f),
    \qquad
    u_i=A(z)^{-1}f_i .
\]
For a test function $v_r$, define the exact weak row and label, with
$j=(i,r)$,
\[
    g_{j,k}:=\langle A_ku_i,v_r\rangle,
    \qquad
    b_j:=\langle f_i-A_0u_i,v_r\rangle .
\]
Then, by the PDE,
\[
    b_j=g_j^\top z .
\]
Thus, for a fixed task $z$, the whole prompt induces the exact linear system
\[
    G_z z=b_z,
    \qquad
    G_z\in\R^{M\times K},
    \qquad
    M=mR.
\]
Conditionally on $z$, every row $g_j$ is Gaussian because it is a linear image
of the Gaussian forcing $f_i$:
\[
    g_j\mid z
    \sim
    N\!\left(0,H_r(z)\right),
    \qquad
    H_r(z)
    =
    B_rA(z)^{-1}\Sigma_fA(z)^{-\top}B_r^\top,
\]
where $(B_ru)_k=\langle A_ku,v_r\rangle$.
The exact task-conditional prompt law is therefore
\[
    P_z(dG,db)
    =
    P_z(dG)\,
    \delta_{Gz}(db),
\]
where $P_z(dG)$ is the Gaussian law determined by the matrices
$\{H_r(z)\}_{r=1}^R$ and the repeated forcing draws.

\begin{proposition}[exact correspondence]
\label{prop:exactcorrespondence}
The PDE ICL training distribution is exactly a low-rank task distribution over
prompts:
\[
    z^\mu\sim N(0,\Sigma_z),
    \qquad
    \calC^\mu\sim P_{z^\mu},
    \qquad
    b^\mu=G^\mu z^\mu .
\]
The task is $z^\mu$; the prompt is the whole matrix--vector pair
$(G^\mu,b^\mu)$ generated from the many PDE pairs
$\{(f_i^\mu,u_i^\mu)\}_{i=1}^{m}$.
\end{proposition}

\begin{proof}
For every $i,r$,
\[
    b_{(i,r)}
    =
    \langle f_i-A_0u_i,v_r\rangle
    =
    \left\langle
    \sum_{k=1}^K z_k A_ku_i,
    v_r
    \right\rangle
    =
    \sum_{k=1}^K z_k g_{(i,r),k}.
\]
Stacking all $mR$ rows gives $b=Gz$.  Since $f_i$ is Gaussian and
$u_i=A(z)^{-1}f_i$, $G$ is conditionally Gaussian given $z$, with covariance
given above.  This proves the prompt-level correspondence.
\end{proof}

\begin{remark}
With $f\sim GP$, the design law $P_z(dG)$ depends on the task through
$A(z)^{-1}$.  Therefore the exact GP-forcing model is not the paper's
task-independent iid Gaussian design unless $H_r(z)$ is independent of $z$ and
isotropic after normalization.  The paper is recovered exactly in that special
case, or in the alternative data-generation convention where $u_i$ is sampled
from a task-independent GP and $f_i=A(z)u_i$.
\end{remark}

\section{Task Diversity and GP-Induced Difficulty}

The exact task diversity object is the empirical covariance of the distinct
pretraining tasks:
\[
    S_z
    :=
    \frac1{M_0}\sum_{\mu=1}^{M_0}z^\mu(z^\mu)^\top .
\]
In the population model, $S_z$ concentrates around
$\Sigma_z=\E[zz^\top]$.  Its rank, or more precisely its spectrum, is the exact
analogue of the paper's low-rank task matrix.  Thus
\[
    r_z=\operatorname{rank}(\Sigma_z),
    \qquad
    \rho_z=\frac{r_z}{K},
    \qquad
    \kappa_z=\frac{M_0}{K}.
\]
The task-diversity transition is
\[
    M_0\simeq r_z
    \qquad
    \text{or equivalently}
    \qquad
    \kappa_z\simeq\rho_z .
\]

The GP forcing does not create new tasks.  It determines how well each task is
observed through the prompt.  The exact observability, or prompt difficulty, is
the collection of conditional row covariances
\[
    \mathcal H
    :=
    \{H_r(z):z\in\operatorname{supp}(\Sigma_z),\ r=1,\ldots,R\},
    \qquad
    H_r(z)=B_rA(z)^{-1}\Sigma_fA(z)^{-\top}B_r^\top .
\]
For a scalar summary one may inspect the spectrum of
\[
    H(z):=\frac1R\sum_{r=1}^{R}H_r(z),
\]
especially restricted to the active task subspace
$\operatorname{range}(\Sigma_z)$.  Small eigenvalues of this restricted matrix
mean that some coefficient directions are almost invisible to the GP prompts.
This is difficulty, not diversity.

\section{Replica Error Reused from the Paper}

Let $\Pi\in\{\mathrm{TM},\mathrm{IDG},\mathrm{ODG}\}$ denote task
memorization, in-distribution generalization, or out-of-distribution
generalization.  In the notation of the paper, the exact task correspondence is
\[
    D\mapsto K,
    \qquad
    r\mapsto r_z,
    \qquad
    \rho=r/D\mapsto \rho_z=r_z/K,
    \qquad
    \kappa=M_0/D\mapsto \kappa_z=M_0/K .
\]
For the exact GP prompt law $P_z(dG,db)$, the same replica construction applies
to a Hamiltonian averaged over the joint law $(z,G,b)$ with $b=Gz$.  Denote the
resulting scalar errors by
$E_{\rm TM}^{\rm rep,GP}$, $E_{\rm IDG}^{\rm rep,GP}$, and
$E_{\rm ODG}^{\rm rep,GP}$.

If, in addition, the prompt design is task-independent and isotropic,
$H_r(z)=I_K/K$, the exact GP errors reduce to the paper's closed-form
low-rank-design errors
\[
    E_{\rm TM}^{\rm rep}
    =
    1-2m+q+\frac{\bar q}{\widetilde\alpha},
    \qquad
    E_{\rm IDG}^{\rm rep}
    =
    1-2m_0+q_0+\frac{\bar q}{\widetilde\alpha},
    \qquad
    E_{\rm ODG}^{\rm rep}
    =
    1-2\bar m+\bar q+\frac{\bar q}{\widetilde\alpha},
\]
where $(q,m,q_0,m_0,\bar q,\bar m)$ solve the same self-consistent equations
as in Appendix E of \cite{takanami2026lowrankicl}.  In the noiseless
large-pretraining limit $\alpha\to\infty$, $\sigma=0$, $\kappa_z\gamma>1$,
the paper's closed form becomes
\[
E_{\rm TM}^{\rm rep}
=
\frac{\min(\rho_z,\kappa_z)}{\widetilde\alpha}
\left(
1+\frac{1-\min(\rho_z,\kappa_z)}{\kappa_z\gamma-1}
\right),
\]
and
\[
E_{\rm IDG}^{\rm rep}
=
\begin{cases}
E_{\rm TM}^{\rm rep}, & \rho_z<\kappa_z,\\[2mm]
1-\dfrac{\kappa_z}{\rho_z}
-\dfrac{\kappa_z(\kappa_z-\rho_z)}
        {\rho_z(\kappa_z\gamma-1)}
+E_{\rm TM}^{\rm rep}, & \rho_z\ge \kappa_z .
\end{cases}
\]
This is the phase transition: below $\rho_z=\kappa_z$ the learned map exploits
the low-rank task structure; above it, performance is bottlenecked by task
diversity.

\section{Encoder as Low-Rank ICL}

The encoder is trained from PDE prompts to output weak rows and labels
$(\wh G,\wh b)$.  Its relevant generalization metric is not coefficient error
but weak consistency:
\[
    \calE_{\rm enc}
    :=
    \frac1M
    \E\|\wh Gz-\wh b\|_2^2 .
\]
For each fixed task $z$, all $M=mR$ rows in the prompt obey
$b_j=g_j^\top z$.  If the encoder is implemented by a one-layer linear
attention model, its scalar prediction for a queried weak row $q$ has the form
\[
    \wh b(q)
    =
    \tr(W_{\rm enc}H_q^\top),
    \qquad
    H_q
    =
    \frac{K}{L_{\rm enc}}
    q\sum_{\ell\le L_{\rm enc}} b_\ell x_\ell^\top ,
\]
with $x_\ell$ the representation of the $\ell$-th weak row.  The task variable
used by the replica calculation is $z$; the full prompt generated from
$P_z(dG,db)$ is the context for that task.

\begin{proposition}[encoder residual bound]
\label{prop:encoder}
Suppose the encoder uses linear attention and is pretrained on the exact
GP-induced prompt law of Proposition~\ref{prop:exactcorrespondence}.  Under
the replica saddle-point prediction for this exact law,
\[
    \boxed{
    \calE_{\rm enc}
    =
    \frac1M\E\|\wh Gz-\wh b\|_2^2
    \le
    E_{\Pi,\rm enc}^{\rm rep,GP}.
    }
\]
When $H_r(z)=I_K/K$ exactly, this is the paper's
$E_{\Pi,\rm enc}^{\rm rep}$ with $D=K$ and $r=r_z$.
\end{proposition}

\begin{proof}
By Proposition~\ref{prop:exactcorrespondence}, each prompt is sampled by first
drawing a single task $z$, then drawing the whole matrix--vector pair
$(G,b)\sim P_z$ with $b=Gz$.  Thus every scalar residual component is a
prediction error for the same task $z$:
\[
    \wh b_j-\wh g_j^\top z .
\]
The replicated partition function for the encoder is therefore the paper's
linear-attention partition function with $D=K$, but with the example
distribution inside each context replaced by the exact conditional law
$P_z(dG,db)$.  The saddle-point prediction of this replicated system is, by
definition, $E_{\Pi,\rm enc}^{\rm rep,GP}$.  Averaging the $M$ scalar residuals
inside a prompt gives the displayed bound.  If $H_r(z)=I_K/K$, the conditional
law is exactly the iid Gaussian design of the paper.
\end{proof}

\section{Decoder as a Learned Preconditioner}

Now set the decoder attention to be linear, without softmax.  Given an encoded
linear system $(\wh G,\wh b)$, define the empirical cross-correlation
\[
    s
    :=
    \frac{K}{M}\wh G^\top\wh b .
\]
For a query vector $q\in\R^K$, the paper's linear-attention feature is
\[
    H_q
    =
    q s^\top .
\]
The decoder prediction is
\[
    \wh y(q)
    =
    \tr(W_{\rm dec}H_q^\top)
    =
    q^\top W_{\rm dec}s .
\]
Thus the decoder's recovered latent vector is
\[
    \boxed{
    \wh z=W_{\rm dec}s
    =
    W_{\rm dec}
    \left(\frac{K}{M}\wh G^\top\wh b\right).
    }
\]
Hence $W_{\rm dec}$ is a learned preconditioner.  In the paper's exact
task-independent Gaussian design, the population normal matrix is the identity:
\[
    \E\left[\frac{K}{M}G^\top G\mid z\right]=I_K.
\]
For the GP-forced PDE prompt law, the exact population normal matrix is instead
\[
    \E\left[\frac{K}{M}G^\top G\mid z\right]
    =
    \frac{K}{R}\sum_{r=1}^{R}H_r(z),
\]
so the learned preconditioner is adapted to this task-dependent geometry.  In
the paper's constant-design case, the replica calculation gives
\[
    W_{\rm dec}^{\star}
    \overset{d}{=}
    \left(
        \sqrt{\wh\chi}\,T+
        \sqrt{\wh{\bar\chi}}\,R+
        \wh m S
    \right)
    \left(
        \wh q S+(\lambda+\wh{\bar q})I_K
    \right)^{-1},
\]
where
\[
    S=\frac1{M_0}\sum_{\mu=1}^{M_0}w^\mu(w^\mu)^\top
\]
is the task covariance and the hatted parameters are the conjugate order
parameters of the paper.  This formula shows explicitly that linear attention
learns a spectral preconditioner adapted to the low-rank task covariance $S$.

\begin{assumption}[bounded learned preconditioner]
\label{ass:boundedprec}
On the high-probability replica event,
\[
    \|W_{\rm dec}^{\star}\|_{\op}\le B_{\rm dec}.
\]
\end{assumption}

\begin{proposition}[decoder latent bound]
\label{prop:decoder}
Assume the decoder is linear attention and the encoded rows are sampled from
the exact GP prompt law of Proposition~\ref{prop:exactcorrespondence}.  Let
\[
    e:=\wh b-\wh Gz
\]
be the encoder residual.  Define the exact GP row-transfer constant
\[
    C_{\rm GP}
    :=
    \frac{
    \E\left[
        \frac{K}{M^2}\|\wh G^\top e\|_2^2
    \right]}
    {
    \E\left[
        \frac1M\|e\|_2^2
    \right]
    },
\]
with the convention $C_{\rm GP}=0$ if the denominator is zero.  Under the
replica saddle-point prediction for the exact GP prompt law and
Assumption~\ref{ass:boundedprec},
\[
    \boxed{
    \frac1K\E\|\wh z-z\|_2^2
    \le
    E_{\Pi,\rm dec}^{\rm rep,GP}
    +
    B_{\rm dec}^2 C_{\rm GP}
    \calE_{\rm enc}.
    }
\]
Combining with Proposition~\ref{prop:encoder},
\[
    \boxed{
    \frac1K\E\|\wh z-z\|_2^2
    \le
    E_{\Pi,\rm dec}^{\rm rep,GP}
    +
    B_{\rm dec}^2 C_{\rm GP}
    E_{\Pi,\rm enc}^{\rm rep,GP}.
    }
\]
If the design is exactly isotropic iid with row covariance $I_K/K$ and the
residual is conditionally independent of the rows, then
$C_{\rm GP}=K/M=1/\widetilde\alpha_{\rm dec}$, recovering the paper's prompt
scaling.
\end{proposition}

\begin{proof}
Write the decoder statistic as
\[
    s
    =
    \frac{K}{M}\wh G^\top\wh b
    =
    \frac{K}{M}\wh G^\top\wh Gz
    +
    \frac{K}{M}\wh G^\top e .
\]
The first term is the noiseless prompt statistic analyzed by the exact GP
replica calculation.  Therefore the scalar prediction error
$\E_q[(q^\top W_{\rm dec}s_0-q^\top z)^2]$ equals
$E_{\Pi,\rm dec}^{\rm rep,GP}$, where
$s_0=(K/M)\wh G^\top\wh Gz$.  Since
$q\sim N(0,I_K/K)$,
\[
    \E_q[(q^\top a)^2]=\frac1K\|a\|_2^2,
\]
so the noiseless latent contribution is $E_{\Pi,\rm dec}^{\rm rep,GP}$ after
normalization by $K$.

For the residual term, no isotropic approximation is used:
\[
    \frac1K\E\left\|
    W_{\rm dec}\frac{K}{M}\wh G^\top e
    \right\|_2^2
    \le
    B_{\rm dec}^2
    \E\left[
        \frac{K}{M^2}\|\wh G^\top e\|_2^2
    \right]
    =
    B_{\rm dec}^2 C_{\rm GP}\calE_{\rm enc}.
\]
The two contributions give the claim.
\end{proof}

\section{PDE Solution Bound}

Let the final reconstructed operator be
\[
    \wh A(\wh z)=\wh A_0+\sum_{k=1}^{K}\wh z_k\wh A_k.
\]
Because the dictionary is identifiable only up to a $GL_K$ gauge, assume we
have fixed the best-conditioned gauge and write
\[
    \epsilon_{\rm dict}
    :=
    \left\|
    \wh A_0-A_0
    \right\|_{\op}
    +
    \|z\|_2
    \left(
    \sum_{k=1}^{K}\|\wh A_k-A_k\|_{\op}^2
    \right)^{1/2},
\]
and
\[
    B_\Psi^2:=\sum_{k=1}^{K}\|\wh A_k\|_{\op}^2 .
\]
Then
\[
    \|\wh A(\wh z)-A(z)\|_{\op}
    \le
    \epsilon_{\rm dict}
    +
    B_\Psi\|\wh z-z\|_2 .
\]

\begin{assumption}[ellipticity]
\label{ass:elliptic}
For all admissible $z$, $A(z)$ is symmetric positive definite and
\[
    \lambda_{\min}(A(z))\ge\mu>0.
\]
On the event considered,
\[
    \|\wh A(\wh z)-A(z)\|_{\op}\le\mu/2 .
\]
\end{assumption}

\begin{theorem}[end-to-end bound]
\label{thm:endtoend}
Under Proposition~\ref{prop:exactcorrespondence},
Assumptions~\ref{ass:boundedprec} and \ref{ass:elliptic}, and the replica
saddle-point prediction for the exact GP prompt law,
\[
\boxed{
    \E\frac{\|\wh u_\star-u_\star\|_2^2}{\|u_\star\|_2^2}
    \le
    \frac{8}{\mu^2}
    \E\epsilon_{\rm dict}^2
    +
    \frac{8B_\Psi^2K}{\mu^2}
    \left(
        E_{\Pi,\rm dec}^{\rm rep,GP}
        +
        B_{\rm dec}^2 C_{\rm GP}
        E_{\Pi,\rm enc}^{\rm rep,GP}
    \right).
}
\]
If the final solve uses ridge
$S_\gamma(B)f=(B^\top B+\gamma I)^{-1}B^\top f$, add the usual deterministic
bias term
\[
    2\frac{\gamma^2}{(\mu^2+\gamma)^2}
\]
and replace $2/\mu$ in the proof by the Lipschitz constant of $S_\gamma$ on
the uniformly elliptic operator class.
\end{theorem}

\begin{proof}
For the exact solve,
\[
    \wh u_\star-u_\star
    =
    \wh A(\wh z)^{-1}
    \bigl(A(z)-\wh A(\wh z)\bigr)u_\star .
\]
Since $\lambda_{\min}(\wh A(\wh z))\ge\mu/2$,
\[
    \frac{\|\wh u_\star-u_\star\|_2}{\|u_\star\|_2}
    \le
    \frac{2}{\mu}\|\wh A(\wh z)-A(z)\|_{\op}.
\]
Using
\[
    \|\wh A(\wh z)-A(z)\|_{\op}^2
    \le
    2\epsilon_{\rm dict}^2
    +
    2B_\Psi^2\|\wh z-z\|_2^2
\]
and Proposition~\ref{prop:decoder} proves the theorem.
\end{proof}

\section{Training Dynamics of the Encoder and Decoder}

To fit exactly into the linear-attention training-dynamics framework, remove
the auxiliary losses $L_u,L_a,L_A$ during this analysis.  Train the encoder
with the prompt consistency loss and train the decoder with the latent
least-squares loss:
\[
    L_{\rm enc}
    =
    \E\left[
        \frac1M\|\wh G_\theta z-\wh b_\theta\|_2^2
    \right],
    \qquad
    L_{\rm dec}
    =
    \frac1K\E\|\wh z_\theta-z\|_2^2 .
\]
The extra physical losses are useful experimentally, but they obscure the
exact correspondence with linear regression ICL.

\subsection{Collapsed Effective Dynamics}

Let $P_{\rm enc}(t)$ and $P_{\rm dec}(t)$ denote the effective linear maps
implemented by the trained linear-attention blocks after collapsing the value
and key-query factors.  This is the analogue of the effective matrix in the
linear-attention dynamics papers.

For a context matrix $G_c$ with labels $b_c=G_cz$, define the context
correlation statistic
\[
    s_c:=\frac{K}{M_c}G_c^\top b_c
    =
    \frac{K}{M_c}G_c^\top G_cz .
\]
For an encoder query weak row $q$, the linear-attention encoder predicts
\[
    \wh b(q,t)=q^\top P_{\rm enc}(t)s_c .
\]
The exact encoder training loss is therefore
\[
    L_{\rm enc}(P)
    =
    \E\left[
        \bigl(q^\top z-q^\top Ps_c\bigr)^2
    \right].
\]
Let
\[
    H_q(z):=\E[q q^\top\mid z],
    \qquad
    C^{\rm enc}_{zs}:=\E[H_q(z)zs_c^\top],
    \qquad
    \mathcal C^{\rm enc}_{ss}(P):=\E[H_q(z)P s_cs_c^\top].
\]
Then the exact gradient-flow dynamics of the collapsed encoder map is
\[
    \boxed{
    \tau_{\rm enc}\dot P_{\rm enc}
    =
    C^{\rm enc}_{zs}
    -
    \mathcal C^{\rm enc}_{ss}(P_{\rm enc}).
    }
\]
If the query row covariance is task-independent, $H_q(z)=H_q$, this becomes
the closed linear matrix ODE
\[
    \tau_{\rm enc}\dot P_{\rm enc}
    =
    H_qC_{zs}-H_qP_{\rm enc}C_{ss},
    \qquad
    C_{zs}:=\E[zs_c^\top],
    \quad
    C_{ss}:=\E[s_cs_c^\top].
\]
When $H_q$ is invertible, the fixed point is
\[
    P_{\rm enc}^{\star}=C_{zs}C_{ss}^{-1}.
\]
Thus encoder training learns the preconditioner that maps the context
correlation statistic $s_c$ back to the task $z$; consistency
$\wh Gz-\wh b\approx0$ is obtained because the predicted row label is
$q^\top P_{\rm enc}^{\star}s_c\approx q^\top z$.

For the decoder, the statistic is
\[
    s:=\frac{K}{M}\wh G^\top\wh b .
\]
The decoder always has the least-squares form
\[
    \boxed{
    \wh z(t)=P_{\rm dec}(t)s.
    }
\]
Training changes only the preconditioner $P_{\rm dec}(t)$.  The exact
collapsed decoder loss is
\[
    L_{\rm dec}(P)=\frac1K\E\|z-Ps\|_2^2,
\]
and gradient flow gives
\[
    \boxed{
    \tau_{\rm dec}\dot P_{\rm dec}
    =
    C^{\rm dec}_{zs}-P_{\rm dec}C^{\rm dec}_{ss},
    \qquad
    C^{\rm dec}_{zs}:=\E[zs^\top],
    \quad
    C^{\rm dec}_{ss}:=\E[ss^\top].
    }
\]
Hence
\[
    P_{\rm dec}^{\star}
    =
    C^{\rm dec}_{zs}
    (C^{\rm dec}_{ss})^{-1},
\]
and
\[
    P_{\rm dec}(t)
    =
    P_{\rm dec}^{\star}
    +
    \bigl(P_{\rm dec}(0)-P_{\rm dec}^{\star}\bigr)
    \exp\!\left(-C^{\rm dec}_{ss}t/\tau_{\rm dec}\right)
\]
whenever $C^{\rm dec}_{ss}$ is constant and diagonalizable.  In an eigenbasis
of $C^{\rm dec}_{ss}$, each mode obeys
\[
    p_a(t)
    =
    p_a^\star
    +
    \bigl(p_a(0)-p_a^\star\bigr)
    e^{-\lambda_a t/\tau_{\rm dec}}.
\]
Large-eigenvalue directions of the prompt statistic are learned first; weakly
observed GP/PDE directions are learned slowly or not at all.

\subsection{Factorized Linear-Attention Dynamics}

The collapsed ODE above describes the algorithm learned by the effective map.
The actual linear-attention parameters are factorized.  This factorization is
what creates the abrupt or staged dynamics observed in the training-dynamics
papers \cite{zhang2025trainingdynamics,mainali2025exactdynamics}.

For merged key-query linear attention, Zhang et al. \cite{zhang2025trainingdynamics}
show that the model is
equivalent to a two-layer linear network on a cubic feature map and that
training has two fixed points: the zero map and the global ICL solution.  From
small initialization, the loss has a long plateau followed by one abrupt drop.
In our notation, that abrupt drop is the moment when
$P_{\rm enc}(t)$ or $P_{\rm dec}(t)$ rapidly moves from near zero to the
preconditioning map $C_{zs}C_{ss}^{-1}$.

For separate key and query, the same paper shows exponentially many fixed
points and saddle-to-saddle dynamics.  The model learns principal components
sequentially.  Translating to the PDE decoder, during the $(m+1)$-st plateau,
\[
    \wh z(t)
    \approx
    \sum_{a=1}^{m}
    p_a^\star e_ae_a^\top s,
\]
where $e_a$ are eigenvectors of the prompt-statistic covariance.  This is
principal-component least squares: the decoder is still applying a
least-squares-style normal-equation statistic, but its preconditioner has rank
$m$.  After all active modes are learned, the decoder becomes the full learned
least-squares preconditioner.

This staged spectral picture matches exact linear-transformer dynamics:
timescale separation is controlled by the data covariance
\cite{mainali2025exactdynamics}.  It also matches softmax-attention analyses:
anisotropic covariates lead to learned preconditioned gradient descent
\cite{he2025demystified}.

For GP-forced PDE prompts, the ordering of the stages is governed by the
spectrum of the exact observable normal statistic
\[
    \Gamma(z):=\frac{K}{M}G_z^\top G_z,
    \qquad
    \bar\Gamma:=\E[\Gamma(z)].
\]
Smooth GP kernels suppress high-frequency forcing modes, which shrinks the
corresponding eigenvalues of $\bar\Gamma$ and delays those coefficient
directions in training.  This is the exact sense in which the GP kernel changes
training difficulty: it changes the spectrum of the statistics that drive the
preconditioner ODE.

\subsection{Training-Dynamics Summary}

The encoder and decoder have the same dynamical skeleton:
\[
    \text{context correlations }s
    \quad\longrightarrow\quad
    \text{learned preconditioner }P(t)
    \quad\longrightarrow\quad
    \text{task estimate }P(t)s .
\]
The encoder uses this skeleton to make the prompt system consistent,
$\wh Gz-\wh b\to0$.  The decoder uses the same skeleton to solve the normal
equations.  At every time, the decoder has the form
\[
    \wh z(t)=P_{\rm dec}(t)\frac{K}{M}\wh G^\top\wh b.
\]
So the algorithmic primitive is always least squares; training dynamics only
determines which spectral preconditioner has been learned, and how many
observable task directions have entered that preconditioner.

\section{Exact MP Training Curves in the Isotropic Design Limit}

This section gives the fully explicit training curves in the regime where the
weak-row design is exactly isotropic Gaussian on the active task subspace.  Let
$B\in\R^{K\times r}$ have orthonormal columns, $r=r_z$, and write
$z=B\zeta$ with
\[
    \zeta\sim N\!\left(0,\frac{K}{r}I_r\right).
\]
Assume
\[
    G B=\frac1{\sqrt K}X,
    \qquad
    X_{ja}\overset{\rm iid}{\sim}N(0,1),
    \qquad
    M=mR.
\]
The normalized active Gram matrix is
\[
    \Gamma
    :=
    \frac{K}{M}B^\top G^\top G B
    =
    \frac1M X^\top X
    \in\R^{r\times r}.
\]
The decoder statistic restricted to the active subspace is exactly
\[
    s_B:=B^\top\frac{K}{M}G^\top b
    =
    \Gamma\zeta .
\]

\begin{theorem}[MP spectrum and exact preconditioner curve]
\label{thm:mpcurve}
Let
\[
    \widetilde\alpha:=\frac{M}{r},
    \qquad
    c:=\frac{r}{M}=\widetilde\alpha^{-1}.
\]
As $M,r\to\infty$ with $r/M\to c$, the empirical spectral distribution of
$\Gamma$ converges almost surely to the Marchenko--Pastur law
\[
    \mu_{\rm MP,c}
    =
    (1-\widetilde\alpha)_+\delta_0
    +
    \frac{\widetilde\alpha}{2\pi\lambda}
    \sqrt{(\lambda_+-\lambda)(\lambda-\lambda_-)}
    \mathbf 1_{[\lambda_-,\lambda_+]}(\lambda)\,d\lambda,
\]
where
\[
    \lambda_\pm=(1\pm\sqrt c)^2
    =
    \left(1\pm\frac1{\sqrt{\widetilde\alpha}}\right)^2 .
\]
Equivalently, the nonzero squared singular values of the raw active design
$GB$ are $(M/K)\lambda$ with $\lambda\sim\mu_{\rm MP,c}$.
Train the collapsed decoder preconditioner with
\[
    L(P)=\frac1K\E_\zeta\|\zeta-P\Gamma\zeta\|_2^2
    =
    \frac1r\operatorname{tr}(I-P\Gamma)(I-P\Gamma)^\top
\]
and gradient flow
\[
    \tau\dot P=\Gamma-P\Gamma^2,
    \qquad
    P(0)=0.
\]
Then, for every finite $r,M$,
\[
    P(t)=\Gamma^\dagger\bigl(I-e^{-\Gamma^2t/\tau}\bigr),
    \qquad
    P(t)\Gamma=\Pi_{\operatorname{range}\Gamma}
    \bigl(I-e^{-\Gamma^2t/\tau}\bigr).
\]
Consequently the exact finite-dimensional training curve is
\[
    \mathcal E_r(t)
    :=
    \frac1K\E_\zeta\|\zeta-P(t)\Gamma\zeta\|_2^2
    =
    \frac1r\operatorname{tr}e^{-2\Gamma^2t/\tau}.
\]
In the MP limit,
\[
    \boxed{
    \mathcal E(t)
    =
    \int e^{-2\lambda^2t/\tau}\,\mu_{\rm MP,c}(d\lambda).
    }
\]
\end{theorem}

\begin{proof}
Since $GB=K^{-1/2}X$,
\[
    \Gamma=\frac{K}{M}B^\top G^\top GB=\frac1M X^\top X,
\]
so the Marchenko--Pastur theorem gives the stated limit.  The loss identity
follows from $\E[\zeta\zeta^\top]=(K/r)I_r$:
\[
    \frac1K\E\|\zeta-P\Gamma\zeta\|_2^2
    =
    \frac1r\operatorname{tr}(I-P\Gamma)(I-P\Gamma)^\top .
\]
Differentiating this quadratic loss gives
\[
    \tau\dot P=\Gamma-P\Gamma^2.
\]
Diagonalize $\Gamma=U\operatorname{diag}(\lambda_a)U^\top$.  Since $P(0)=0$,
the solution remains diagonal in this eigenbasis.  Each eigenmode satisfies
\[
    \tau\dot p_a=\lambda_a-p_a\lambda_a^2,
    \qquad
    p_a(0)=0.
\]
Thus, if $\lambda_a>0$,
\[
    p_a(t)=\lambda_a^{-1}(1-e^{-\lambda_a^2t/\tau}),
    \qquad
    p_a(t)\lambda_a=1-e^{-\lambda_a^2t/\tau},
\]
while if $\lambda_a=0$, $p_a(t)\lambda_a=0$.  Therefore the residual
eigenvalue is $1-p_a(t)\lambda_a=e^{-\lambda_a^2t/\tau}$, including zero
modes.  Averaging the squared residual eigenvalues gives the finite curve; the
MP limit follows by weak convergence of the empirical spectral measure.
\end{proof}

\begin{corollary}[spectrum of the learned normal-equation map]
\label{cor:pgspectrum}
The empirical spectrum of $P(t)\Gamma$ is the pushforward of
$\mu_{\rm MP,c}$ by
\[
    \phi_t(\lambda)=
    1-e^{-\lambda^2t/\tau}.
\]
Equivalently, away from atoms, with
\[
    \lambda(y)=
    \sqrt{-\frac{\tau}{t}\log(1-y)},
    \qquad
    y\in \phi_t([\lambda_-,\lambda_+]),
\]
the density is
\[
    \rho_{P\Gamma,t}(y)
    =
    \rho_{\rm MP,c}(\lambda(y))
    \frac{\tau}{2t\,\lambda(y)(1-y)}.
\]
As $t\to\infty$, $P(t)\Gamma$ converges to the projector onto
$\operatorname{range}\Gamma$.  Hence the limiting error is
\[
    \lim_{t\to\infty}\mathcal E(t)
    =
    (1-\widetilde\alpha)_+ .
\]
\end{corollary}

\begin{corollary}[replica-predicted HB and Chebyshev hyperparameters]
\label{cor:replicahbhyperparameters}
Assume $\widetilde\alpha>1$, $t>0$, and consider the ridge-free active
subspace, so $\lambda_->0$ and $P(t)\succ0$.  Define the replica-predicted
support of the preconditioned normal operator by
\[
    a_-(t,c)=\phi_t(\lambda_-),
    \qquad
    a_+(t,c)=\phi_t(\lambda_+).
\]
Then the distribution-free minimax Heavy--Ball hyperparameters are predicted
without a validation sweep:
\begin{align}
    \alpha_{\rm HB}^{\rm rep}(t,c)
    &=\frac{4}{
        \bigl(\sqrt{a_+(t,c)}+\sqrt{a_-(t,c)}\bigr)^2},
        \label{eq:rep-hb-alpha}\\
    \beta_{\rm HB}^{\rm rep}(t,c)
    &=\left(
        \frac{\sqrt{a_+(t,c)}-\sqrt{a_-(t,c)}}
             {\sqrt{a_+(t,c)}+\sqrt{a_-(t,c)}}
      \right)^2.
        \label{eq:rep-hb-beta}
\end{align}
The same endpoints determine the complete exact Chebyshev coefficient
schedule.

More generally, let $p_L^{\rm HB}(a;\alpha,\beta)$ be the depth-$L$ HB
residual polynomial.  For an isotropic active coefficient vector, the
normalized $\Gamma$-energy error converges almost surely to
\begin{equation}
    \mathcal R_L^{\rm HB}(\alpha,\beta;t,c)
    =
    \frac{
      \displaystyle\int
      \lambda\,
      p_L^{\rm HB}(\phi_t(\lambda);\alpha,\beta)^2
      \,\mu_{\rm MP,c}(d\lambda)
    }{
      \displaystyle\int\lambda\,\mu_{\rm MP,c}(d\lambda)
    }.
    \label{eq:rep-hb-distributional-risk}
\end{equation}
Consequently the distribution-aware prediction is the deterministic
two-scalar problem
\begin{equation}
    (\alpha_L^\star,\beta_L^\star)
    \in\operatorname*{argmin}_{
      0\leq\beta<1,
      \ \alpha a_+(t,c)<2(1+\beta)}
    \mathcal R_L^{\rm HB}(\alpha,\beta;t,c).
    \label{eq:rep-hb-distributional-optimum}
\end{equation}
If $T_s(L)$ is the hardware cost of controller $s$, its predicted depth and
controller are obtained by minimizing the corresponding replica risk plus
$\tau_{\rm hw}T_s(L)$.  Thus $c$, $t/\tau$, and the hardware tradeoff replace
an empirical hyperparameter grid in this isotropic model.
\end{corollary}

\begin{proof}
Corollary~\ref{cor:pgspectrum} gives the support endpoints of $P(t)\Gamma$.
The classical HB minimax formulas on that interval yield
\eqref{eq:rep-hb-alpha}--\eqref{eq:rep-hb-beta}.  Since $P(t)$ is a spectral
function of $\Gamma$, the matrices commute.  Starting from zero, the error is
$p_L^{\rm HB}(P(t)\Gamma)\zeta$.  Isotropy of $\zeta$ turns its expected
$\Gamma$-energy into the normalized trace
\[
    \frac{\tr\!\left[
      \Gamma p_L^{\rm HB}(P(t)\Gamma)^2\right]}{\tr\Gamma}.
\]
Apply the Marchenko--Pastur convergence and the continuous mapping theorem to
obtain \eqref{eq:rep-hb-distributional-risk}; minimizing it gives
\eqref{eq:rep-hb-distributional-optimum}.
\end{proof}

\begin{corollary}[replica-predicted cluster-aware polynomial]
\label{cor:replica-moment-polynomial}
Let \(\gamma_q\) be the positive limiting \(H\)-energy spectral measure
produced by the encoder--head replica order parameters \(q\), and fix a
decoder depth \(L\).  If its moments
\[
 m_k(q)=\int\lambda^k\,\gamma_q(d\lambda),
 \qquad 0\leq k\leq2L,
\]
are finite and \(\mathsf H_L(q)=(m_{i+j}(q))_{i,j=1}^L\) is nonsingular,
then the replica-predicted optimal right-hand-side-independent residual
polynomial is
\[
 p_{L,q}^\star(\lambda)
 =1+\sum_{j=1}^Lc_j(q)\lambda^j,
 \qquad
 \mathsf H_L(q)c(q)=-(m_1(q),\ldots,m_L(q))^\top .
 \label{eq:pde-replica-moment-polynomial}
\]
It minimizes the limiting energy risk over every fixed degree-\(L\)
polynomial controller, and therefore weakly improves on the HB and
support-only Chebyshev predictions at the same depth.  For clustered laws the
improvement can be strict even when the support endpoints and condition
number are unchanged.
\end{corollary}
\begin{proof}
For \(p(\lambda)=1+\sum_{j=1}^Lc_j\lambda^j\), the limiting risk is
\[
 m_0+2\sum_{j=1}^Lc_jm_j
 +\sum_{i,j=1}^Lc_ic_jm_{i+j}.
\]
Differentiate this strictly convex quadratic.  HB and Chebyshev are feasible
members of the same polynomial class.
\end{proof}

This corollary is the principled version of a Chebyshev-weight MLP.  A neural
head may predict promptwise invariant cluster nodes and masses, or the
moments above, but the small Gram solve and polynomial evaluation remain
hard-coded.  In finite precision the Gram system should be expressed in a
shifted Chebyshev basis.  Predicting only \((\lambda_-,\lambda_+)\) recovers
support minimax Chebyshev but discards precisely the multiplicities that
allow a fixed polynomial to approach PCG on clustered PDE spectra.

The spectral measure and the stability interval must be parameterized
separately.  The implemented controller predicts ordered positive nodes,
positive normalized masses, and a positive expansion of a Ritz reference;
the expansion is capped by the matrix-free global upper certificate.  If an
actual eigenvalue exceeds the Clenshaw endpoint \(\widehat L\), then its
shifted coordinate \(x=2\lambda/\widehat L-1\) is larger than one and

\[
 T_j(x)=\cosh\!\left(j\operatorname{arcosh}x\right).
\]

Thus even a vanishingly rare coverage miss can dominate an arithmetic-mean
risk exponentially in depth.  This is an exact polynomial obstruction, not
an approximation-theorem defect of the MLP.  A deterministic result must use
the certified endpoint; a learned tighter endpoint requires an explicit tail
assumption or a residual fallback.

The finite-dimensional audit
\texttt{audit\_pde\_moment\_chebyshev\_learning.py} tests the corollary with
only the spectral-measure MLP trainable.  At $K=32,M=128,S=4,L=8,r=2$,
three seeds and 4096 fresh tasks per seed, the correlated elliptic law gives
mean relative $H$-risk $9.00\times10^{-4}$.  This improves on
oracle-interval Richardson ($3.78\times10^{-1}$), HB
($8.65\times10^{-1}$), and Chebyshev ($7.79\times10^{-2}$), showing that
the learned multiplicities are genuine useful statistics.  It does not
approach same-preconditioner PCG--8 ($1.81\times10^{-8}$) or head-free
PCG--8 ($2.19\times10^{-6}$).  On the well-conditioned elliptic law the
learned risk is $7.44\times10^{-7}$, while the exact-interval stationary
methods are already at the float32 floor.  In both laws the spectral formula
and realized Clenshaw risk agree to plotting precision.  Therefore the
replica moment law is predictive for the implemented controller, but it does
not imply dominance over adaptive Krylov methods.

\begin{remark}[exact input to the deformed PDE law]
For the GP-forced PDE prompt law, one must not simply count the $M=mR$ weak
rows as iid samples.  Conditional on $z$, one forcing produces
$u\sim\mathcal N(0,C_z)$ with
\[
    C_z=A(z)^{-1}\Sigma_fA(z)^{-\top},
\]
and its $R$ weak rows share this same $u$.  If $b_{rk}^\top$ is row $k$ of
$B_r$ and
\[
    S_{k\ell}
    :=\frac12\sum_{r=1}^R
    (b_{rk}b_{r\ell}^\top+b_{r\ell}b_{rk}^\top),
\]
then one forcing contributes the random normal-matrix entry
$Y_{k\ell}=u^\top S_{k\ell}u$.  Gaussian Wick calculus gives the exact
conditional first two moments
\begin{align}
    \E[Y_{k\ell}\mid z]
    &=\tr(S_{k\ell}C_z),\\
    \operatorname{Cov}(Y_{k\ell},Y_{pq}\mid z)
    &=2\tr(S_{k\ell}C_zS_{pq}C_z).
    \label{eq:pde-normal-wick}
\end{align}
The prompt normal matrix is a sum of $m$ independent copies of this
\emph{matrix-valued quadratic form}.  Equation \eqref{eq:pde-normal-wick},
not an iid-row MP ansatz, is therefore the correct finite-$K$ input to a CLT,
matrix Dyson equation, or operator-valued replica closure.  It also recovers
the exact population matrix
\[
    \E\left[\frac K{mR}G^\top G\mid z\right]
    =\frac KR\sum_{r=1}^R H_r(z).
\]

An idealized head that routes individual eigendirections has a scalar,
commuting closure.  With \(S/K\to\sigma\), limiting query law \(\pi\), and
bounded invariant score \(\ell(q,\lambda;\nu)\), define
\begin{align*}
    Z(q)&=\int e^{\ell(q,t;\nu)}\nu(dt),\\
    g(\lambda;\nu)
    &=1-\exp\left[-\sigma\int
      \frac{e^{\ell(q,\lambda;\nu)}}{Z(q)}\pi(dq)\right],\\
    T_\theta(\lambda;\nu)&=(1-g(\lambda;\nu))\lambda+g(\lambda;\nu).
\end{align*}
The effective law is exactly the pushforward
$\nu_{\rm eff}=(T_\theta)_\#\nu$.  Therefore
\eqref{eq:rep-hb-distributional-risk}--
\eqref{eq:rep-hb-distributional-optimum} apply with $\phi_t$ replaced by
$T_\theta$ and jointly predict the head parameters, slot density, HB pair,
and depth for that commuting oracle.  This is not the exact closure of the
implemented block-power matrix-free head unless its routed subspace becomes
invariant.

For the implemented head, let \(Y_\theta\in\mathbb R^{K\times S}\) be the
softmax-routed block before QR,
\(s_H=\operatorname{tr}(H)/\overline L\), \(A=H/s_H\), and \(r\) the fixed number of
block-power refinements.  Assuming \(A^rY_\theta\) has full column rank, set
\begin{align}
 K_r&=Y_\theta^\top A^{2r}Y_\theta,
 &J_r&=Y_\theta^\top A^{2r+1}Y_\theta,\\
 U&=A^rY_\theta K_r^{-1/2},
 &C&=K_r^{-1/2}J_rK_r^{-1/2},\\
 L_r&=Y_\theta^\top A^{2r+2}Y_\theta,
 &R=AU-UC,&
 R^\top R&=K_r^{-1/2}L_rK_r^{-1/2}-C^2.
 \label{eq:pde-mf-ritz-interface}
\end{align}
If \(c_\star=\lambda_{\min}(C)\), let
\(M=c_\star C^{-1}\),
\(d_\perp=\overline L-\operatorname{tr}C\), and
\[
 \gamma=\|RM^{1/2}\|_{\mathrm F},\qquad
 \widehat L_{\rm post}
 =\frac{c_\star+d_\perp+
 \sqrt{(c_\star-d_\perp)^2+4\gamma^2}}{2},qquad
 \zeta=\widehat L_{\rm post}/\overline L.
\]
The actual preconditioner is exactly
\begin{equation}
 B_\theta=(s_H\zeta)^{-1}\left[
 I-UU^\top+c_\star UC^{-1}U^\top\right].
 \label{eq:pde-mf-preconditioner}
\end{equation}
Thus the neural component learns \(Y_\theta\), whereas all covariance
contractions, the \(S\times S\) inverse, and the solver recurrence are exact.
Because \(U\) need not be an invariant subspace of \(H\), the replica closure
requires the limits of \((K_r,J_r,L_r)\) and the cross measures of \(U\) with the
teacher, right-hand side, and physical decoder metric.  A scalar eigenvalue
pushforward alone is insufficient; no layer-indexed order parameter is needed
for HB or Chebyshev once these prompt-only quantities are fixed.

There is also a sharp scaling boundary.  If
\(A_\theta=B_\theta^{1/2}HB_\theta^{1/2}\), then
\[
 \operatorname{rank}(A_\theta-H/(s_H\zeta))\leq2S,
 \qquad
 \sup_x|F_{A_\theta}(x)-F_{H/(s_H\zeta)}(x)|\leq 2S/K.
\]
A fixed-slot head therefore cannot alter an isotropic normalized bulk risk by
order one.  It can alter the low-rank PDE risk by order one when the physical
or teacher subspace has nonvanishing overlap with
\(\operatorname{span}(U,HU)\).  This overlap, rather than a claim of universal
spectral improvement, is the first-principles source of benefit against an
unconditioned solver.  A bulk change requires \(S=\Theta(K)\), concentrating
attention, or an equivalent proportional-rank mechanism.

The finite-dimensional correlated-PDE audit at
$K=32,M=128,S=4,r=2,L=8$ verifies that the scalar certificate order parameter
$\zeta$ is not cosmetic.  Across three seeds, replacing trace-only scaling by
the post-deflation block bound reduces trained HB relative $H$-risk from
$6.88\times10^{-1}$ to $7.34\times10^{-5}$.  The same prompt head followed by
PCG--8 reaches $1.81\times10^{-8}$, versus $4.74\times10^{-8}$ for head-free
PCG--20 at equal sequential scalar-HVP work.  Thus the block moments and
$\zeta$ predict a real amortized geometry benefit, while the PCG boundary
remains necessary: HB--8 is still less accurate than head-free PCG--8.

Gaussian forcing is used only for the explicit Wick covariance above; the
finite-dimensional predictor requires only an SPD normal matrix, finite
moments, and the stated nondegeneracy of the routed block.
\end{remark}

\begin{proposition}[prompt-free finite-$K$ hyperparameter prediction]
\label{prop:prompt-free-pde-predictor}
Fix $K,R$ and condition on a task $z$.  Let
\[
    \overline H_m(z)
    :=m\,\E[Y\mid z]+\lambda M,
    \qquad
    \overline c_m(z)
    :=m\,\E[G_i^\top b_i\mid z],
\]
where both expectations are the exact Gaussian contractions above and $M$ is
the covariant ridge metric.  For probe \(r\), write \(g_r(u,z)^\top\) for its
weak row, \(v_r=g_r/\|g_r\|\), and let \(\varphi_z(g_r)\) be the invariant row
feature after replacing prompt centering and scaling by their conditional
population values.  For slot \(s\), define the population-routed direction
\[
 \overline y_{\theta,s}(z)
 =
 \frac{\displaystyle\sum_{r=1}^R
  \E\!\left[
   e^{\ell_{\theta,s}(\varphi_z(g_r))}
   v_r\mid z\right]}
 {\displaystyle\sum_{r=1}^R
  \E\!\left[
   e^{\ell_{\theta,s}(\varphi_z(g_r))}\mid z\right]},
 \qquad
 \overline Y_\theta(z)
 =[\overline y_{\theta,1},\ldots,\overline y_{\theta,S}].
 \label{eq:population-routed-block}
\]
Let \((H_m,c_m,Y_{\theta,m})\) be the normal matrix, right-hand side, and
softmax-routed block from an actual length-$m$ prompt.  If the score/value
moments in \eqref{eq:population-routed-block} are finite and its population
feature normalization is nondegenerate, then
\[
    m^{-1}(H_m-\overline H_m(z))\longrightarrow0,
    \qquad
    m^{-1}(c_m-\overline c_m(z))\longrightarrow0,
    \qquad
    Y_{\theta,m}-\overline Y_\theta(z)\longrightarrow0
\]
almost surely.  If the population normal matrices remain uniformly SPD, the
invariant one-head score is Lipschitz on the resulting compact spectral set,
and the routed block remains uniformly full rank, the effective spectral
law, every fixed-depth HB or Chebyshev risk, and
the corresponding physical linearized risk computed from
$(\overline H_m,\overline c_m,\overline Y_\theta)$ converge to their
prompt-law counterparts.
Their conditional fluctuations are $O_{\mathbb P}(m^{-1/2})$.

If, additionally, the latency-regularized population objective has a unique
interior minimizer with nonsingular Hessian, its prediction
\[
 (\theta^\star,\alpha^\star,\beta^\star,L^\star)
 \in\arg\min_{\theta,\alpha,\beta,L}
 \left\{
 \E_z\mathcal R_L^{\rm HB}
   (\overline H_m(z),\overline c_m(z),\overline Y_\theta(z);
    \theta,\alpha,\beta)
 +\tau_{\rm hw}T_{\rm HB}(L,\theta)
 \right\}
\]
is consistent for the finite-prompt optimum.  The analogous statement holds
for the two Chebyshev endpoints and their exact recurrence.
\end{proposition}

\begin{proof}
The prompt is a sum over $m$ independent forcings, not over $mR$ independent
rows.  The strong law applied to the finite-dimensional matrix-valued
quadratic form $Y$ and its cross moment proves the first two almost-sure
limits.
Apply the same law to the numerator, denominator, and empirical feature
normalization of each softmax slot.  The denominator is strictly positive,
so the continuous-mapping theorem gives
\eqref{eq:population-routed-block}.  Gaussian quadratic forms have finite
moments of every order, and the assumed score/value moments permit the
multivariate central limit theorem and delta method; the conditional
fluctuations are therefore $O_{\mathbb P}(m^{-1/2})$.
Uniform SPD, uniform rank of the routed block, continuity of QR and of the
small Ritz functional calculus, and the Lipschitz spectral head pass this
convergence through \eqref{eq:pde-mf-ritz-interface}--
\eqref{eq:pde-mf-preconditioner} and every fixed-degree residual polynomial.
The physical statement follows by the same argument for the bounded decoder
Jacobian.  Standard argmin stability
under the stated local identifiability condition gives consistency of the
joint minimizer.
\end{proof}

This proposition separates two meanings of the phrase few assumptions.  At
finite $K$, the conditional-population predictor needs no MP, isotropy,
eigenvector delocalization, or replica ansatz.  Eliminating the expectation
over the task prior $z$ as well---so that only dimension ratios
remain---does require the stronger deformed random-matrix closure.  No method
can identify a distribution-optimal controller from dimension ratios alone
when two task laws with the same ratios have different spectral weights.

\begin{corollary}[early and late time laws]
\label{cor:earlylate}
At small time,
\[
    \mathcal E(t)
    =
    1-\frac{2t}{\tau}\int\lambda^2\,\mu_{\rm MP,c}(d\lambda)
    +O(t^2)
    =
    1-\frac{2t}{\tau}(1+c)+O(t^2).
\]
If $\widetilde\alpha>1$, then $\lambda_->0$ and
\[
    \mathcal E(t)
    \asymp
    \exp(-2\lambda_-^2t/\tau)\,t^{-3/2}
    \qquad
    (t\to\infty).
\]
If $\widetilde\alpha=1$, then $\lambda_-=0$ and
\[
    \mathcal E(t)\asymp t^{-1/4}.
\]
If $\widetilde\alpha<1$, the curve converges to the unobservable-dimensional
plateau $1-\widetilde\alpha$.
\end{corollary}

\begin{remark}
The same theorem applies to the encoder after replacing $M$ by the number of
weak rows in the encoder context and $\tau$ by $\tau_{\rm enc}$.  Its curve is
the consistency curve for $\wh Gz-\wh b$.  For the exact GP-forced PDE law, the
MP law is replaced by the spectral law of the deformed sample covariance with
population covariance $H(z)$.  Thus the formulas above are exact for the
isotropic reduction and become deformed-MP formulas for the full GP geometry.
\end{remark}

\section{What Has Been Reused}

The imported statistical-mechanics step is the one from
\cite{takanami2026lowrankicl}: the trained linear-attention matrix is replaced
by a replica quadratic surrogate, whose solution is a spectral shrinkage or
preconditioning map.  In the exact GP-forcing PDE model, the surrogate is
averaged over the task-conditional prompt law $P_z$, not over a
task-independent iid Gaussian design.  The paper's closed forms are recovered
when $H_r(z)=I_K/K$ exactly.

Everything specific to the PDE prompt law is exact:
\[
    f\sim GP
    \Rightarrow
    u=A(z)^{-1}f\ \text{Gaussian conditional on }z
    \Rightarrow
    g=\bigl(\langle A_ku,v\rangle\bigr)_{k=1}^{K}\ \text{Gaussian}
    \Rightarrow
    b=g^\top z .
\]
Thus the encoder residual is a full-prompt low-rank ICL prediction error for
the task $z$, and the decoder is a learned-preconditioner problem applied to
the normal-equation statistic $(K/M)G^\top b$.

\begin{remark}[practical reading]
The theorem predicts that improving the encoder residual helps the decoder only
through the exact factor $C_{\rm GP}$.  In the isotropic iid design case,
$C_{\rm GP}=K/M$; for GP-forced PDE prompts, $C_{\rm GP}$ is determined by the
actual weak-row geometry.  The hard task-diversity transition is still governed
by whether the number of distinct pretraining tasks $M_0$ is large enough
relative to the latent rank $r_z$.
\end{remark}

\begin{thebibliography}{1}
\bibitem{takanami2026lowrankicl}
Kaito Takanami, Takashi Takahashi, and Yoshiyuki Kabashima.
\newblock Learning Linear Regression with Low-Rank Tasks In-Context.
\newblock \emph{AISTATS 2026 / arXiv:2510.04548}, 2026.
\newblock \url{https://arxiv.org/pdf/2510.04548}.

\bibitem{zhang2025trainingdynamics}
Yedi Zhang, Aaditya K. Singh, Peter E. Latham, and Andrew Saxe.
\newblock Training Dynamics of In-Context Learning in Linear Attention.
\newblock \emph{ICML 2025 / arXiv:2501.16265}, 2025.
\newblock \url{https://arxiv.org/abs/2501.16265}.

\bibitem{mainali2025exactdynamics}
Nischal Mainali and Lucas Teixeira.
\newblock Exact Learning Dynamics of In-Context Learning in Linear Transformers
and Its Application to Non-Linear Transformers.
\newblock \emph{arXiv:2504.12916}, 2025.
\newblock \url{https://arxiv.org/abs/2504.12916}.

\bibitem{he2025demystified}
Jianliang He, Xintian Pan, Siyu Chen, and Zhuoran Yang.
\newblock In-Context Linear Regression Demystified: Training Dynamics and
Mechanistic Interpretability of Multi-Head Softmax Attention.
\newblock \emph{ICML 2025 / arXiv:2503.12734}, 2025.
\newblock \url{https://arxiv.org/abs/2503.12734}.
\end{thebibliography}

\end{document}
